\chapter{Installation and setup}\label{ch:setup}

\section{Requirements}

\begin{longtable}{@{}P{0.24\linewidth} P{0.70\linewidth}@{}}
\caption{Requirements.}\label{tab:requirements}\\
\toprule
Item & Requirement \\
\midrule
\endfirsthead
\toprule
Item & Requirement \\
\midrule
\endhead
OpenFOAM & ESI v2412 with cfMesh (\cmd{cartesianMesh}), and its Open MPI. \\
Tools & \cmd{bash}, \cmd{rsync}, \cmd{awk}, GNU \cmd{sed}, \cmd{xz}, \cmd{/usr/bin/time}, \cmd{jinja2-cli}. \\
Python 3 & Standard library; \cmd{matplotlib} for \file{cdFlight.py}. \\
\cmd{perf} & Optional: counts the CPU cycles and the energy of each run. \\
Memory & Proportional to the number of cells. Measure \ent{maxRSSkB} of a run on the coarsest mesh (\cref{sec:runinfo}) and scale it with the number of cells before you run a finer mesh. \cmd{cartesianMesh} needs the memory of the full rocket also for a half mesh. \\
\bottomrule
\end{longtable}

\section{Prepare the shell}

Do this step in each new shell:
\begin{shell}
source /usr/lib/openfoam/openfoam2412/etc/bashrc
\end{shell}

\begin{caution}
On WSL, put the mesh and run directories on the Linux file system, not in
\file{/mnt/c}: the input and output of OpenFOAM become several times slower.
\end{caution}

\section{Prepare the computer for a study}\label{sec:prepare-computer}

The cost of a run is a measurement only if the computer runs at a known, constant
state (\cref{sec:measurement}). Before a study:

\begin{steps}
\item Fix the CPU frequency policy: the \val{performance} governor or platform
  profile. A power-saving policy lowers the clock and changes the cost.
\item Let \cmd{perf} count events: \ent{kernel.perf_event_paranoid} $\leq 1$ for the
  cycles, $\leq 0$ also for the energy. Make the setting permanent in
  \file{/etc/sysctl.d/}.
\item Prevent suspension during the study (power settings, or \cmd{systemd-inhibit}).
\item Give the memory and the CPU sufficient cooling. Above its temperature limit,
  memory refreshes more often and loses bandwidth; a CPU lowers its clock.
\item Run one case on each CPU socket, with nothing heavy beside it. Runs on one
  socket share the memory bandwidth.
\end{steps}

\section{Build this guide}\label{sec:build-guide}

The guide needs LuaLaTeX, BibTeX and the TeX Gyre fonts. The diagrams need Node.js,
the npm package \texttt{pele} and Google Chrome.

\begin{steps}
\item Go to \file{docs/guide/}.
\item If you changed a diagram in \file{figures/src/*.mmd}, render the diagrams:
\begin{shell}
npm install
make figures
\end{shell}
  Pele renders the Mermaid source to SVG; headless Chrome prints it to a vector PDF,
  because Pele sets the colours with CSS variables.
\item Compile with \cmd{make}. Commit the source, \file{figures/pdf/} and the PDF.
\end{steps}
